% Folder 91, batch 03 — archivists' pages 41–60.
% Pass: Opus 5.5 (claude-opus-5-5), 2026-10-03 — first pass, unchecked against the pages by a human.
%
% Source: archives/batches/91/batch-03.pdf (cover sheet excluded), read from
% the 700-dpi tiles in archives/tiles/91/p41 … p60. The hand is a fast
% draft throughout: the formulas mostly carry, the connective prose often
% does not, and the apparatus says so.
\documentclass[11pt,a4paper]{article}
\input{../preamble/grothendieck}

\folder{91}
\batch{3}
\pages{41}{60}
\foldertitle{Autour de Néron / Greenberg-Néron. Foncteurs Hom (méthodes non-projectives) : notes manuscrites (s.d.), lettre (1967).}
\dating{[à partir de 1964-vers 1970]}
\watermark{Édition de démonstration}

\begin{document}

\page{41}
\note{la page s'ouvre au milieu d'un argument commencé avant ce lot (sur la représentabilité de $f_{*}(X/T)$, cf. p. 43 et suivantes).}

d'où l'hypothèse \struck{\ill{} $(f^{-1}(U_i))$ \ill{} $T$} ; pour tout $s\in S$, il existe un $i(s)$ tel que
$\varphi(T_s)\subset U_{i(s)}$, il \uncertain{en résulte} \ill{} (précision) que l'on peut trouver un voisinage
$V(s)$ de $s$ tel que $\varphi(T|V(s))\subset U_i$. \uncertain{Cela} \ill{} \uncertain{se traduit} par un
\uncertain{homomorphisme} $V(s)\to H_i$, d'où $V(s)\to H$. \ill{} pas \ill{} \uncertain{veut}. (l'on \ill{}
\ill{} --- .) Je \uncertain{laisse} \ill{} \uncertain{straight-forward} \ill{}.

\textbf{Théorème.} Les hypothèses \uncertain{sont} \ill{} $T$ fini et loc.\ libre sur $S$
\struck{et $X/S$ quasi-projectif. \ill{} \ill{}}
Les hypothèses \ill{} \ill{} \struck{\ill{} \ill{} \ill{}}
\begin{enumerate}
  \item[(i)] $T\to S$ \ill{} \ill{} \uncertain{fini et loc.\ libre} ;
  \item[(ii)] pour tout $t\in S$, \uncertain{toute partie finie} de $X_s$ \uncertain{au-dessus de} $T_s$, définie sur
  une extension finie de $k(s)$, \uncertain{soit contenue} dans un ouvert affine de $X$.
  \struck{[\ill{} \ill{} \ill{}]}
\end{enumerate}
\marginal{$k(s)$ \quad $k(t)$}
\struck{\ill{} \ill{} \ill{}} \ill{} $f_{*}(X/T)$ existe et est représentable \ill{} \ill{}
$f_{*}(U_i/T)$, où $U_i$ \uncertain{parcourt} \ill{} \uncertain{ouverts affines} de $X$.
\note{les lignes (i)--(ii) sont reprises dans une rédaction très surchargée : une première formulation biffée
est encerclée et reportée par des traits ; l'ordre ci-dessus est celui que les traits indiquent, à titre de
proposition.}

\page{42}
\textbf{Cor. 1.} Les \uncertain{conditions} du \uncertain{Théor.} ci-dessus sont \uncertain{satisfaites} dans
\uncertain{chacun des cas suivants} :
\begin{enumerate}
  \item[(i)] $T/S$ \ill{} \struck{\ill{} \ill{}} \ill{} ;
  \item[(ii)] $S=\bigcup S_i$ \ill{} tels que \uncertain{pour tout} $i$, $X|S_i$ soit \uncertain{contenu dans un}
  \ill{} affine \uncertain{[c'est le cas si} $X$ \uncertain{quasi-proj.\ sur} $S$].
\end{enumerate}
\note{au-dessus de (ii), une ligne interlinéaire biffée et encadrée : \struck{\uncertain{pour chaque} \ill{} \ill{}}.}

\textbf{Cor. 2.} Sous les conditions \uncertain{du Théor.}, soit $X'$ un sous-schéma \ill{} de $X$, \add{posons
$f'=f|X'$}. Alors $f'_{*}(X')$ \uncertain{existe}, et est un ss-schéma \ill{} de $f_{*}(X)$.

\page{43}

\begin{tikzcd}
  & X \arrow[d] \\
  S & T \arrow[l, "f"']
\end{tikzcd}

$f_{*}(X/T)$, \quad $T$ \uncertain{fini sur} $S$.

Si $X/T$ \uncertain{est} affine, $f_{*}(X/T)$ \uncertain{aussi} ;

\quad " \quad projectif, propre \quad ———— \quad pourvu que $f$ \uncertain{séparable} (*) ;

\quad " \quad quasi-projectif, \quad ————.
\note{les deux tirets reprennent « $f_{*}(X/T)$ aussi » ; une accolade réunit les lignes « projectif, propre » et
« quasi-projectif » devant la condition sur $f$.}

$f_{*}$ est exact à gauche, \struck{mais \ill{} \ill{} \ill{}} \add{\uncertain{en particulier} \ill{}}
\[
  f_{*}(T/T)=S,\qquad f_{*}(X\times_T Y/T)=f_{*}(X/T)\times_S f_{*}(Y/T),
\]
\[
  f_{*}(X\times_Z Y/T)=f_{*}(X/T)\times_{f_{*}(Z/T)}f_{*}(Y/T).
\]
$f_{*}$ \uncertain{transforme} immersions \uncertain{en} immersions, \ill{} \ill{} \ill{}. \struck{\ill{}}
\uncertain{immersions} \ill{}, \struck{[\ill{} \ill{}]}
\[
  f_{*}(X/T)=\bigcup_i f_{*}(U_i/T),
\]
\ill{} les $U_i$ \ill{} \ill{} $X$]. \struck{\ill{} \ill{} \ill{}}
[N.B. \uncertain{cette} \ill{} \uncertain{doit pouvoir} \ill{} \ill{} \ill{} \uncertain{démontrer}.]
\add{\ill{} $=$ \uncertain{immersion et immersion} \ill{}}

\uncertain{De même}, $f_{*}$ \uncertain{transforme} \struck{\ill{}} \uncertain{morphismes affines en morphismes affines}
[\ill{} \ill{} \ill{} $\varphi\colon X\to Y$ \uncertain{affine}, \ill{} \ill{} \add{\uncertain{on se ramène au cas où}}
$Y$ \uncertain{est affine} et $X\subset Y\otimes_T Z$, \add{$Z/T$ \uncertain{affine}} \ill{} \ill{},
\ill{} $X=Y\otimes_T Z$ \ill{} \ill{} \ill{}] il \uncertain{transforme} \uncertain{morphismes projectifs en
morphismes projectifs} \uncertain{si} $T/S$ \uncertain{séparable}, et il \uncertain{transforme} \uncertain{morphismes
quasi-projectifs en morphismes quasi-projectifs} \add{\uncertain{si} $X/T\to Y/T$ \ill{}}.
[\uncertain{Même méthode}, \uncertain{en utilisant} \ill{} \ill{} $Y$ \uncertain{est} \ill{}. \struck{\ill{}}
\ill{} \uncertain{proj.} \ill{} \uncertain{de type fini sur} $T$ \uncertain{dans un faisceau} \ill{} sur $Y$ \ill{}
\ill{} \uncertain{d'un} \ill{} de $T$ ; \ill{} \ill{} \ill{}, \ill{} \ill{} \ill{} \uncertain{qu'un} \ill{}
\note{la suite à la p. 44.}

\marginal{(*) $f$ \uncertain{soit pour descente} : \uncertain{si} $Y/S$ \uncertain{est tel que} $f^{*}(Y)/S'$ \uncertain{soit}
propre resp.\ projectif, alors $Y/S$ l'est. [Dans \uncertain{ce} \uncertain{dernier cas}, on \uncertain{suppose}
$T/S$ \uncertain{fini et plat sur} $S$ \uncertain{loc.\ noeth.}, (i) \uncertain{faudrait dans le cas actuel}
\uncertain{pouvoir l'exorciser}, \ill{} \ill{} \uncertain{recouvrement en projectifs types}, \struck{\ill{}}
\uncertain{travailler sur} \ill{} \ill{} \uncertain{en effet}, il \uncertain{faut prouver} \uncertain{ceci} :
\uncertain{si un faisceau loc.\ libre inversible} $\underline{L}$ \uncertain{sur} $X$ \uncertain{est} $T$-\uncertain{ample},
\uncertain{alors} $f_{*}(\underline{L})$ [\uncertain{qui est un faisceau loc.\ libre} \ill{} \uncertain{sur}
$f_{*}(X/T)$ \ill{} \ill{} \ill{}] \uncertain{donne} $\det f_{*}(\underline{L})$ \ill{} $S$-\uncertain{ample}.
\uncertain{Dans le premier cas}, \uncertain{on utilise} \ill{} \uncertain{critère de Chow}, \uncertain{en notant que si}
$X'\to X$ \uncertain{est surjectif}, \uncertain{alors} $f_{*}(X')\to f_{*}(X)$ \uncertain{l'est aussi}
\struck{\ill{}}, \ill{} \ill{} \ill{} \ill{} $S=$ (\ill{}), \ill{} \ill{} \ill{} \uncertain{descente}…
\uncertain{Ici il faut} $S$ \uncertain{loc.\ noeth.}]}

\marginal{\uncertain{Question}. \uncertain{Si} $T/S$ \uncertain{est fidèlement plat} \ill{}, \uncertain{et} $Y/S$ \ill{}
\ill{}, \uncertain{et si} $Y\times_S T/T$ \uncertain{soit propre}, \uncertain{est-il vrai que} $Y/S$ \uncertain{l'est} ?
— \uncertain{Oui} \ill{} \uncertain{dans le cas} $S$ \uncertain{loc.\ noeth.}, \uncertain{car} \ill{} \uncertain{utiliser}
\ill{} \ill{} \uncertain{critère de Chow}.}
\note{les deux paragraphes précédents occupent la colonne de gauche de la page, sous le diagramme ; le premier est l'appel
de l'astérisque. Une note marginale verticale, en partie biffée, longe le diagramme : \ill{}.}

\page{44}
\uncertain{prouver} le \uncertain{dit} \ill{} \uncertain{faisceau descendu}, on \uncertain{trouve} \ill{} ; \uncertain{fonction}
\ill{} \ill{} \uncertain{faisceau} \ill{} \ill{} \uncertain{sur} $X/Y$, \ill{} \uncertain{faisceau} \ill{} \uncertain{sur}
$f_{*}(X/T)/f_{*}(Y/T)$ \uncertain{etc}]. Kif-kif \ill{} \ill{} \ill{}.

\textbf{Proposition.} Soit $g\colon X/T\to Y/T$ un \uncertain{morphisme}, et soit $\underline{L}$ un faisceau
\uncertain{inversible} \uncertain{dans} $X$ \uncertain{ample} \struck{\ill{}} (\uncertain{resp.\ très ample})
\uncertain{relativement à} $g$. \uncertain{Alors} $\det f_{*}(\underline{L})$ \uncertain{est un faisceau ample}
(\uncertain{resp.\ très ample}) \uncertain{pour}
\[
  f_{*}(X/T)\to f_{*}(Y/T).
\]
\marginal{$f_{*}(\ill{})\to f_{*}\mathcal{O}(X/T)$}

\textbf{Démonstration} (2). On \uncertain{peut supposer} $Y$ \uncertain{affine}, \uncertain{l'on} \ill{}
\struck{\ill{} \ill{}} (\uncertain{c'est une notion locale}). (1) \struck{\ill{} \ill{}}
\ill{} « \ill{} » \struck{\ill{}}. \ill{} \uncertain{on se ramène} \ill{} à \uncertain{prouver ceci} : si
$X=\mathbf{P}^{n}_{Y}$, et $\underline{L}$ le \uncertain{faisceau} \ill{}, \ill{} \uncertain{alors}
$f_{*}(\underline{L})/T$ \uncertain{est} \ill{} \uncertain{immersion projective} \ill{} de $f_{*}(\mathbf{P}^{n}_{Y})$
\uncertain{dans} $f_{*}(Y/T)$. Mais \uncertain{utilisant la} \ill{} de $f_{*}$ \ill{}, \ill{} \ill{} \ill{} \ill{}
$Y=T$. \uncertain{Dans ce cas}, \ill{} \ill{} \ill{} \ill{} $S$ \uncertain{affine} $=\operatorname{Spec}(A)$,
\[
  T=\operatorname{Spec}(B),
\]
$B$ \uncertain{étant libre de type fini} (\uncertain{rang} $\nu$) \uncertain{sur} $A$.

\ill{} \ill{} \ill{} \uncertain{à deux vérifications} :

A) \uncertain{Détermination} \ill{} \ill{} \ill{} \uncertain{grassmannienne}, \uncertain{i.e.\ d'une}
\uncertain{immersion} \ill{} \ill{} \ill{} \uncertain{de la droite}.

\marginal{(2) \ill{} \uncertain{remarques} \ill{} \ill{} \ill{}}
\marginal{N.B. \uncertain{Pour} \ill{} \ill{} \ill{} \ill{} (i) $f_{*}(\underline{L})$ \ill{} \ill{} \ill{}
\ill{} (ii) \ill{} \ill{} \ill{} \ill{} \ill{} \ill{}}
\note{la marge gauche porte, d'en bas en haut et dans un cadre irrégulier, une note N.B. à deux points (i), (ii),
presque entièrement illisible ; elle concerne $f_{*}(\underline{L})$. La « Démonstration » est appelée « (2) » ;
son renvoi « (1) » n'a pas de texte lisible.}

\page{45}
B) \uncertain{Étude de la situation suivante} : \ill{} \ill{} \ill{} \ill{}
\ill{} $G$ \ill{} \ill{} \ill{} \ill{} \ill{} \ill{} \uncertain{(question)}
\ill{} \uncertain{faisceau loc.\ libre} \uncertain{de rang} $\nu$ \ill{} \uncertain{d'algèbres}
$\underline{A}$ \uncertain{sur} $G$ (\uncertain{l'image inverse} \ill{} $f(e_T)$)]
\ill{} \uncertain{faisceau loc.\ libre} \uncertain{de} $\underline{A}$ \struck{\ill{}} \ill{} \ill{}
$\underline{F}$, \ill{} \ill{} \ill{} \uncertain{faisceau} \uncertain{de} $\underline{O}_G$\struck{\ill{}}\uncertain{-modules},
\add{\uncertain{tel que} $\underline{F}/\underline{F}_1$ \uncertain{soit}} $\underline{F}_1$, \uncertain{loc.\ libre}
\uncertain{de} \uncertain{rang} $k$. \ill{} \ill{} \ill{} \ill{} \ill{} \ill{} \ill{} \ill{}.

\ill{} \ill{} \ill{} \ill{} $G'$ \uncertain{sur} $G_T$ [\uncertain{exemple} $G\times_S S'$]
\ill{} $\underline{F}'_1/\underline{F}'$ \uncertain{est} \uncertain{loc.\ libre} \ill{}
\struck{$\underline{F}'_1\subset\underline{F}'$ \ill{} \ill{} \ill{}}
\ill{} \uncertain{sur} $\underline{A}'$. \uncertain{On veut} \ill{}
\uncertain{représenter} \ill{} \struck{\ill{}} $G_{*}$ \uncertain{de} $G$ \ill{} \ill{} \ill{} \uncertain{et} \ill{}
\uncertain{universel} : \ill{} \ill{} $G'\to G$ \uncertain{se factorise par} $G_{*}$.

\ill{} la \struck{\ill{}} \add{\uncertain{question}} \ill{} \ill{} \ill{} $G'$ (\ill{}
\uncertain{de la représentabilité} \uncertain{d'un foncteur} !),

\ill{} \ill{} \uncertain{supposons} :
\begin{enumerate}
  \item[(i)] $G$ \uncertain{affine} \uncertain{d'anneau} $A$ ;
  \item[(ii)] $\underline{A}$ \uncertain{libre} \ill{} \ill{} \ill{} $A$-\uncertain{algèbre} $B$ \uncertain{ayant une base}
  $b_1,\dots,b_\nu$ ;
  \item[(iii)] $\underline{F}$ \ill{} \ill{} \ill{} $B^{n}$, \uncertain{d'où} \ill{} \ill{} \uncertain{base}
  $b^{(i)}_{\alpha}$ $\bigl(\begin{smallmatrix}1\le i\le n\\ 1\le\alpha\le\nu\end{smallmatrix}\bigr)$ ;
  \item[(iv)] $\underline{F}_1$ a une \uncertain{base} $(c_\lambda)_{1\le\lambda\le(n-k)\nu}$, \uncertain{qui se}
  \ill{} \ill{} \ill{} \ill{} \ill{} \uncertain{base} \uncertain{d'un} \ill{}
  \uncertain{supplémentaire} $(d_j)_{1\le j\le k\nu}$.
\end{enumerate}
\struck{$\underline{F}=\underline{A}^{n}$, $\underline{F}_1=$}
\note{la formule biffée ci-dessus est barrée d'un zigzag, à droite de « supposons ».}

\uncertain{Les} \ill{} \ill{} \ill{} \uncertain{conditions} \ill{} \uncertain{et coefficients}
\ill{}

\page{46}
\uncertain{pour que} $\underline{F}_1$ \uncertain{soit} \ill{} \ill{} \uncertain{dites}. \struck{\ill{}}

1°) $\underline{F}_1$ \uncertain{stable par} \uncertain{multiplication} par $\underline{B}$,
\uncertain{i.e.}\ \ill{} \ill{} \ill{} $b_i c_\lambda$ \uncertain{s'exprime} \ill{} \ill{}
\ill{} $\varphi_{i\lambda\mu}\in A$), \uncertain{doivent} \uncertain{être} \uncertain{nulles}, \ill{} \ill{} \ill{}
\ill{} \uncertain{équations},
\[
  \varphi_{i\lambda\mu}=0\qquad
  \begin{bmatrix}1\le i\le\nu\\ 1\le\lambda\le(n-k)\nu\\ 1\le\mu\le k\nu\end{bmatrix}
\]
\marginal{\uncertain{Soit} $G_1$ \uncertain{le ss-schéma} \ill{} \ill{} \uncertain{défini} \ill{} \ill{} \ill{}
\ill{} \ill{}}
\note{la note marginale, encerclée, renvoie par un trait à « équations ».}

2°) Le \uncertain{quotient} \struck{\ill{}} $F/F_1=M$ \uncertain{doit} \uncertain{être} \uncertain{projectif} \ill{}
\uncertain{de rang} \ill{}. \ill{} \ill{} \ill{} à la \uncertain{base} $(m_i)_{1\le i\le k\nu}$
\add{\uncertain{des images des} \ill{}} de $M$. \uncertain{Soient} $\xi=(\xi_1,\dots,\xi_k)$ \ill{}
\uncertain{éléments} de $M$. \uncertain{Pour que} \ill{} \ill{} \uncertain{base} \ill{} $B$, il \ill{} \ill{} \ill{}
\ill{} \ill{}
\[
  B^{k}\xrightarrow{\;u_\xi\;}M
\]
\uncertain{défini} \ill{} \uncertain{soit} \uncertain{bijectif}]. \uncertain{Or}, \ill{} \ill{} \ill{} $B^{k}$ \ill{}
$M$ \ill{} \ill{} \ill{}, \ill{} \uncertain{condition} \ill{} \uncertain{que} $\det(u_\xi)$ \uncertain{soit}
\uncertain{inversible}. \ill{} \ill{} \ill{} \ill{} \ill{} \ill{}, \uncertain{pour que} \ill{} \ill{} $x\in\operatorname{Spec}(A)$,
$M_x$ \uncertain{soit libre sur} $B_x$, \uncertain{il faut et il suffit qu'il existe}
$\xi\in(M_x)^{k}$ \uncertain{tel que} $\det(u_\xi)$ \uncertain{soit inversible} \uncertain{dans} $A_x$.
\ill{} \ill{} \ill{} \uncertain{à prendre} \ill{} $\xi\in M^{k}$. \ill{} \ill{} $\xi\in M^{k}$, \uncertain{soit} $U_\xi$
\uncertain{l'ouvert} \ill{} \ill{} \ill{} $\det(u_\xi)$. \uncertain{Soit} $U=\bigcup U_\xi$. (\ill{}
\marginal{\ill{} \ill{} \ill{}, \ill{} \ill{} \uncertain{conditions} \ill{} \ill{} \ill{}}
\note{la page s'arrête sur une parenthèse ouverte ; la suite est p. 47.}

\page{47}
l'\uncertain{ouvert cherché} [N.B. \uncertain{le localisé} \ill{} de $M$ \ill{} \ill{} \ill{} \uncertain{condition}
\uncertain{locale}, \ill{} \uncertain{compatibilité} \ill{} \ill{} \ill{} \uncertain{extensions} \ill{} \ill{}
\uncertain{corps de base} $k(x)$ \ill{}]

\note{deux traits horizontaux séparent ce qui précède de ce qui suit.}

\uncertain{Sous les réserves habituelles}, si $g\colon X/T\to Y/T$ est \uncertain{de type fini}, \uncertain{alors}
$f_{*}(g)$ est \uncertain{de type fini}.
[$S$ \uncertain{loc.\ noeth.}] \ill{} \ill{} \ill{} \ill{} \uncertain{au cas où} $Y=T$,
\struck{$X$ \uncertain{affine}} \add{\ill{}}. \struck{\ill{} \ill{} \ill{} \ill{}}
\struck{$X$} \ill{} \uncertain{localisation} d'un \ill{} \add{\ill{} $U$} \ill{} \ill{} \struck{$T$.} \ill{}
\ill{} \uncertain{hypothèses}, \struck{\ill{}} \ill{} \ill{} $X_s$ \ill{} \ill{} $T_s\to$ \ill{} \ill{} \ill{} \ill{}
\ill{} \ill{}. \ill{}, \uncertain{fini} \ill{} \ill{} \ill{} \ill{} \ill{} \uncertain{affine} \ill{}
\ill{}. \struck{\ill{} \ill{} \ill{}}
\struck{\ill{} \ill{}, \ill{} \ill{}}
\[
  X^{\nu}=\underbrace{X\times_S\cdots\times_S X}_{\nu},
\]
\ill{} \ill{} \ill{} \uncertain{affine} \ill{} $U$ de $X$, \uncertain{soit}
$U^{\nu}=\underbrace{U\times_S\cdots\times_S U}_{\nu}$, $U^{\nu}\to$ \ill{} \ill{}
\ill{} de $X^{\nu}$. \struck{\ill{} \ill{} \ill{} \ill{}} \ill{} $X^{\nu}\to$ \ill{}
\uncertain{utilisation}, il \uncertain{contient} \ill{} \ill{} \uncertain{fini} \ill{}
\uncertain{affines} $U_i$ \uncertain{tels que}
\[
  \bigcup U_i^{\nu}=\bigcup U^{\nu}.
\]
\uncertain{Soit} $(x_1,\dots,x_\nu)$ \struck{\ill{} \ill{} \ill{} \ill{}}
\struck{$U$. \uncertain{Alors}} \ill{} $x\in X$ \ill{} \uncertain{la forme} \ill{} $x_i$.
\uncertain{Mais} \ill{} $x_i$ \ill{} \ill{} \ill{} \ill{} \uncertain{affine} \ill{}
\marginal{\ill{} \ill{} \ill{} \ill{}}
\note{la page se termine en cours de phrase ; la suite est p. 48.}

\page{48}
$U$ \uncertain{ssi} $x\in U^{\nu}$, \struck{\ill{}} \uncertain{soit} \ill{}
$x\in U_i^{\nu}$ \uncertain{pour un} $i$ \uncertain{convenable} i.e.\ \uncertain{les} $x_\alpha\in U_i$.
Cqfd.

\uncertain{Cela est ainsi} \ill{} \ill{} \uncertain{dans le cas où} $X$ \uncertain{est affine sur} $Y$. Mais
\uncertain{alors} \ill{} \ill{} \ill{} \ill{} $X=\operatorname{Spec}(\underline{F})$, \ill{}
\uncertain{faisceau} \uncertain{loc.\ libre de rang} \ill{} \ill{} \ill{} \uncertain{sur} $T$. \uncertain{Alors}
\[
  f_{*}(X/T)=\operatorname{Spec}\bigl(f_{*}(\underline{F})\bigr).
\]

\struck{\ill{} \ill{}}

\textbf{Remarque.} Il \uncertain{n'est pas vrai} \ill{} \ill{} \ill{} \uncertain{que la} \ill{} \uncertain{de}
$f_{*}(X/T)/S$ \uncertain{soit} \ill{} \uncertain{fois la dimension} \uncertain{de} $X/T$.
Ex.\,: $S=\operatorname{Spec}(k)$, $T$ \struck{$=\operatorname{Spec} k[E]/\ill{}$} \uncertain{est} \ill{}
\uncertain{radiciel} \uncertain{sur} $S$, $X=T\times_S T$ i.e.\ $f_{*}(X/T)=\underline{\operatorname{Hom}}_S(T,T)$.
\uncertain{Ce dernier} \ill{} \ill{} \ill{} \uncertain{dimension} \add{$>0$} \struck{\ill{}} \uncertain{sur} $k$.
\uncertain{Exemple} : \uncertain{le} \ill{} \uncertain{des} \uncertain{automorphismes} \ill{} \uncertain{la}
$k$-\uncertain{algèbre} \ill{} \uncertain{est} \ill{} \uncertain{de dim} $>0$.
\marginal{\ill{} \ill{} \ill{} \ill{}}

\page{49}

(A)

\begin{tikzcd}
  f_{*}(X/T) \arrow[r] & f_{*}(Y/T) \\
  S'_1 \arrow[u] \arrow[r, hook] & S' \arrow[u] \arrow[ul, dashed]
\end{tikzcd}

(B)

\begin{tikzcd}[column sep=small, row sep=small]
  & T & & Y & \\
  S & & T' \arrow[ul] \arrow[ur] \arrow[rr, dashed] \arrow[dl] & & X \arrow[ul] \\
  & S' \arrow[ul] & & T'_1 \arrow[ul] \arrow[ur] \arrow[dl] & \\
  & & S'_1 \arrow[ul] & &
\end{tikzcd}

\note{dans (B), le trait entre $T$ et $S$ est appuyé et porte une pointe vers $S$, et celui entre $S'$ et $S$
est sans pointe lisible ; on les a rendus $T\to S$ et $S'\to S$. Les flèches $T'_1\to T'$ et $T'\to S'$ sont
lues comme sur la page, sans en vérifier la cohérence.}

\struck{Proposition} $\{$\uncertain{Théorème} ?$\}$ \uncertain{Soit} $X/T\xrightarrow{g}Y/T$ \uncertain{un}
\uncertain{morphisme} \add{\uncertain{de type fini}} \struck{\uncertain{étale}} \uncertain{de} $T$-\uncertain{préschémas},
$T$ \struck{\ill{}} \struck{\uncertain{sur} $S$ \uncertain{loc.\ noeth.}}
\add{\uncertain{sur} $S$ (\ill{} \uncertain{fini} \ill{}) *}.
\uncertain{Les} \uncertain{hypothèses} \uncertain{que} $f_{*}(X/T)$ \uncertain{et} $f_{*}(Y/T)$ \uncertain{existent} \uncertain{et}
\struck{\ill{} \uncertain{sur} $S$, \ill{} \ill{} $f_{*}(X/T)\to f_{*}(Y/T)$}
\add{\ill{}} $S$ \uncertain{et} $T$ \ill{} \uncertain{loc.\ noeth.}
\begin{enumerate}
  \item[1.] \struck{\uncertain{loc.\ de type fini}} \add{\uncertain{loc.\ de type fini sur}} \ill{} ;
  \item[(i)] \uncertain{Si} \ill{} $T/S$ \struck{\ill{}} \ill{}, \uncertain{si} $X/T\xrightarrow{g}Y/T$ \uncertain{est}
  \uncertain{lisse}, \uncertain{il en est de même de} $f_{*}(g)$ ;
  \item[(ii)] \uncertain{Si} $T/S$ \struck{\ill{}} \add{\uncertain{de type fini}}, \uncertain{si} $S$ \ill{}
  \uncertain{stable}, $f_{*}(g)$ \uncertain{est étale}.
\end{enumerate}

\textbf{Démonstration.} (i) Il \uncertain{faut prouver} \ill{} \ill{} \ill{} \ill{} \uncertain{le}
\uncertain{diagramme} (A), \ill{} $S'_1\to S'$ \uncertain{affine}, \add{\ill{} \uncertain{de type fini} $/S$} $S'_1$
\ill{} \struck{\ill{} $S$} \uncertain{ss-schéma} \ill{} $=$ \ill{} \ill{}.
\ill{} \ill{} : \ill{} \uncertain{condition} \ill{} \ill{} \ill{} \uncertain{que}
\ill{} \uncertain{diagramme} \ill{} (B), (\ill{} \ill{}), \ill{} $T$ \uncertain{affine}, $T'$ \ill{} \ill{}
\ill{}, \uncertain{et} $T'_1\to T'$ \ill{} \ill{} \uncertain{immersion} \ill{} \ill{} \ill{}. \ill{} $T'_1$ \ill{}
\ill{} \ill{} $T'$ \ill{} $=$ \uncertain{ss-schéma} \ill{} \ill{}
\marginal{\uncertain{Cela suppose} $S$ \ill{}}
\note{la phrase se poursuit à la p. 50.}

\page{50}
(ii) \uncertain{Ce} \uncertain{raisonnement} \uncertain{vaut} \ill{} \ill{} \uncertain{supposer} $T/S$
\uncertain{étale}, \uncertain{si} $g$ \uncertain{est} \uncertain{étale}, \struck{\ill{}} \uncertain{étale}, $f_{*}(g)$
\uncertain{est} \uncertain{non ramifié}. \uncertain{Pour} \uncertain{prouver} \uncertain{qu'il} \uncertain{est} \uncertain{étale},
il \uncertain{suffit} \uncertain{de} \uncertain{prouver} (\uncertain{par} \ill{} \uncertain{libre}, \ill{}
\struck{\ill{}} \add{\uncertain{fini}} \uncertain{ensuite}) \struck{\ill{} \ill{} \ill{}}.
\uncertain{On se ramène} \uncertain{au cas où} $S=\operatorname{Spec}(k)$, $k$ \uncertain{alg.\ clos}, \uncertain{et}
\uncertain{prouver} \ill{} \uncertain{pour} \uncertain{si} $k'\to$ \ill{}
\marginal{\uncertain{Dans le} \uncertain{cas} \uncertain{non fini}}

\begin{tikzcd}
  & f_{*}(X) \arrow[d, leftrightarrow] \\
  \operatorname{Spec}(k') \arrow[r] & f_{*}(Y)
\end{tikzcd}

\note{à gauche du diagramme, « $k$ » au-dessus de $\operatorname{Spec}(k')$ et « $T$ » au-dessus de la colonne
$f_{*}$ ; la flèche verticale porte des pointes aux deux bouts, ou une pointe surchargée.}

\uncertain{extension} \uncertain{nilpotente} \uncertain{de} $k$, \uncertain{alors} il \ill{}
\ill{} … \uncertain{qui} \uncertain{se} \uncertain{relève} \uncertain{en} $k'$-\uncertain{point}
\ill{} $f_{*}(X)$ \uncertain{au-dessus} \uncertain{d'un} $k'$-\uncertain{point} \ill{} $f_{*}(Y)$, i.e.\ \uncertain{que}
\ill{} \ill{} \ill{} \uncertain{qui} \ill{} \ill{} \uncertain{fini}
\ill{} \uncertain{section}. \ill{} $X'/T'$ \uncertain{qui} \uncertain{relève} \uncertain{une} \uncertain{section}
\uncertain{donnée} \uncertain{de} $Y'/T'$.
[\ill{} \ill{} \ill{} \ill{}. \uncertain{La solution est} \ill{}]
\uncertain{Or} $T'$ \ill{} \uncertain{nilpotent} \ill{} [\uncertain{de type fini sur} $k'$]
\ill{} \uncertain{donc} \ill{} \uncertain{les} \ill{} \uncertain{composantes}. \uncertain{Donc} \ill{} \ill{}
\uncertain{que} \ill{} $T'$ \uncertain{connexe}. \uncertain{Alors} \ill{} \ill{} \uncertain{comme} \ill{}
\uncertain{lorsque} \ill{} \uncertain{réduit} \ill{} \ill{} \uncertain{l'image} \uncertain{d'un} \uncertain{pt}.
\uncertain{Cela} \uncertain{démontre} \ill{} \ill{}.

\textbf{\uncertain{Corollaire} 1.} \uncertain{Soit} $\nu$ \uncertain{le} \uncertain{maximum} \uncertain{des}
\uncertain{rangs} \struck{\ill{}} \uncertain{des} \uncertain{anneaux} \ill{} \ill{} \uncertain{locaux}
\ill{} $T_s$ ($s\in S$), \uncertain{et soit} \ill{} $n$ \uncertain{le} \uncertain{maximum} \uncertain{atteint} \uncertain{par}
\ill{} $X/Y$. \uncertain{Alors} \uncertain{le} \uncertain{maximum}
\note{la page s'arrête en cours de phrase ; la suite est p. 51.}

\page{51}
\uncertain{atteint} \uncertain{par} \uncertain{les} \uncertain{fibres} \uncertain{de}
$f_{*}(X/T)\to f_{*}(Y/T)$ \uncertain{est} $\le n\nu$.
\struck{\uncertain{Corollaire}} \struck{[\uncertain{Si} $T/S$ \uncertain{est} \ill{} \uncertain{fini}]}

\struck{\uncertain{Corollaire} 1. \uncertain{Soit} \ill{} \ill{} \ill{}}

\struck{\ill{} \ill{} \ill{}}

\struck{\uncertain{Soit} $d=\sup_{y\in Y}\dim g^{-1}(y)$, \uncertain{et}}

\struck{$d'=\sup_{\bar y\in f_{*}(X/T)}\dim f_{*}(g)^{-1}(\bar y)$.}

\struck{\uncertain{Soit} $n=\deg(T/S)$.}

\struck{\uncertain{Alors} \ill{} \ill{} $d'=nd$.}

\struck{\uncertain{On a égalité si}}

\struck{\uncertain{si} \ill{} \ill{} \uncertain{les fibres}}

\struck{\uncertain{On a égalité si}}

\struck{$\deg_{\text{sép}}(T/S,s)$}

\struck{\uncertain{est indépendant de}}

\struck{$s$ (\uncertain{égal à} $\nu$)\,]}

\struck{\uncertain{si les} \uncertain{fibres} \uncertain{de} $f$ \ill{}}

\struck{\uncertain{toutes} \uncertain{les} \uncertain{dimensions} \ill{} \ill{}}

\struck{\uncertain{de} $X$, \uncertain{toutes} \uncertain{égales},}

\struck{\uncertain{il en est de même de}}

\struck{\uncertain{celles} \ill{} $=$.}
\note{tout ce passage, du premier « Corollaire 1 » au « $=$ » qui le termine, est biffé par de longs traits et
enfermé dans un contour ; on l'a rendu biffé ligne à ligne.}

\textbf{\uncertain{Corollaire} 2.} \uncertain{Supposons} $T/S$ \uncertain{fini} \uncertain{et de} \uncertain{degré}
$n$ \uncertain{constant}. \uncertain{Supposons} \ill{} \ill{} \ill{}, \uncertain{les} \uncertain{conditions} (i).
\uncertain{Soit} \struck{\ill{} \ill{}} $s\in S$, \uncertain{soit} $D(s)$ \uncertain{l'ens.\ des}
\uncertain{dimensions} \uncertain{des} \uncertain{fibres} \uncertain{composantes} \uncertain{des}
\uncertain{fibres} \uncertain{de} $X_s\to Y_s$ ; \struck{\ill{}} \uncertain{soit} $D'(s)$ \uncertain{l'ens.}
\ill{} \ill{} \uncertain{pour} $f_{*}(X/T)_s\to f_{*}(Y/T)_s$.
\uncertain{Alors} $D(s)=n\,D'(s)$.\note{ainsi sur la page ; la conclusion qui suit (dimension $nd$ pour $f_{*}(g)$) demanderait plutôt $D'(s)=n\,D(s)$.} \uncertain{En particulier}, \ill{} \ill{} \uncertain{les}
\uncertain{composantes} \uncertain{de} \uncertain{toutes} \uncertain{les} \uncertain{fibres} \uncertain{de} $g$ \ill{}
$=$ \uncertain{dimension} $d$, \uncertain{toutes} \uncertain{les} \uncertain{composantes} \uncertain{de}
\uncertain{toutes} \uncertain{les} \uncertain{fibres} \uncertain{de} $f_{*}(g)$ \ill{} $=$ \uncertain{dimension} $nd$.
\note{dans la marge gauche, en regard de l'énoncé : $X_s\to Y_s$, puis $k(s)$ et $T_s$.}

\page{52}
\textbf{\uncertain{Démonstration}.} \uncertain{Les} \uncertain{deux} \uncertain{membres} \uncertain{supposés}
\ill{} $S=\operatorname{Spec}(k)$, \uncertain{De plus}, \uncertain{on} \uncertain{peut} \uncertain{supposer} $k$
alg.\ clos. \uncertain{Ensuite}, \uncertain{d'après}
\[
  \Bigl(\coprod f_i\Bigr)_{*}(X/T)=\prod_i f_{i*}(X_i/T_i),
\]
\uncertain{on} \uncertain{se} \uncertain{ramène} \uncertain{au} \uncertain{cas} \uncertain{où} $T$ \uncertain{est}
\uncertain{local}, i.e.\ $T\to S$ \uncertain{radiciel} \uncertain{et} \uncertain{surjectif}. \uncertain{On}
\uncertain{peut} \uncertain{aussi} \uncertain{supposer} $X$ \uncertain{et} $Y$ \uncertain{affines}, \uncertain{et}
$X\xrightarrow{g}Y$ \uncertain{se} \uncertain{factorisant} \uncertain{en}
\[
  X\xrightarrow{g_1}Y[t_1,\dots,t_d]\xrightarrow{g_2}Y,
\]
\struck{$X\to X$} \uncertain{où} $g_1$ \uncertain{est} \uncertain{étale} \uncertain{et} $g_2$ \uncertain{la}
\uncertain{projection} \uncertain{canonique}.
\struck{\uncertain{D'où} $f_{*}(g)=f_{*}(g_2)\,f_{*}(g_1)$}
\struck{\uncertain{et} $f_{*}(g_1)$ \uncertain{étale},} \uncertain{et} \uncertain{il} \uncertain{suffit}
\uncertain{de} \uncertain{prouver} \uncertain{que} \uncertain{les} \uncertain{fibres} \uncertain{de}
$f_{*}(g)$ \uncertain{sont} \uncertain{de} \uncertain{dimension} $nd$.
\uncertain{Comme} $f_{*}(g)=f_{*}(g_2)\,f_{*}(g_1)$, \uncertain{et} $f_{*}(g_1)$ \uncertain{est}
\uncertain{étale} (\uncertain{et} \ill{} \ill{} \uncertain{surjectif}), \uncertain{on} \uncertain{est}
\uncertain{ramené} \uncertain{au} \uncertain{cas} \uncertain{où}
\[
  X=Y[t_1,\dots,t_d]=Y\times_T\overline{T}[t_1,\dots,t_d]
\]
\struck{\uncertain{i.e.} $2\cdot T[t]$}
\uncertain{On} \uncertain{est} \uncertain{ramené} (\uncertain{grâce} \uncertain{à} \uncertain{(i)}) :
\uncertain{prouver} \uncertain{que} \uncertain{les} \uncertain{fibres} \uncertain{de}
\[
  f_{*}\bigl(T[t_1,\dots,t_d]/T\bigr)
\]
\uncertain{sont} \uncertain{de} \uncertain{dim} $nd$, \uncertain{ce} \uncertain{qui} \uncertain{est} \uncertain{trivial}….
\note{le $\overline{T}$ (une barre au-dessus d'un $T$ biffé) est une lecture douteuse.}

\page{53}

\begin{tikzcd}[column sep=small]
  X_s \arrow[r] & Y_s & & X'_s \arrow[r] & Y'_s \\
  k(s)=k \quad T_s & & k' & T'_s \arrow[u] \arrow[ur] &
\end{tikzcd}

\note{sous $T_s$, un premier essai biffé : \struck{$X'_s$ \ill{} $T'_s$}.}

\marginal{Nb \uncertain{de} \uncertain{comp.\ connexes} \uncertain{géométriques} \uncertain{de} $T_s$ : $\nu_s$}
\[
  \begin{cases}
    \nu=\sup_s\nu_s\\
    n=\deg T/S\\
    d=\dim.\ \text{fibres de } g\\
    \lambda=\deg X/Y
  \end{cases}
\]
\note{ces quatre définitions, en haut à droite, sont réunies par une accolade et un chevron.}

\textbf{\uncertain{Corollaire} 1 bis.} \uncertain{Sous} \uncertain{les} \uncertain{conditions} \uncertain{du}
\uncertain{corollaire} 1,
\struck{\uncertain{Supposons} \uncertain{de} \uncertain{plus} $T/S$ \uncertain{fini},}
\struck{\uncertain{Alors}} \uncertain{et} \uncertain{supposons} \uncertain{que}
\uncertain{le} \uncertain{morphisme} \uncertain{étale} \struck{\ill{}} $g\colon X\to Y$ \uncertain{soit} \uncertain{propre}
\uncertain{et} \uncertain{de} \uncertain{degré} $\lambda$. \uncertain{Alors}
$f_{*}(g)_s\colon f_{*}(X/T)_s\to f_{*}(Y/T)_s$ \uncertain{est} \uncertain{propre} \uncertain{et}
\uncertain{de} \uncertain{degré} $\lambda^{\nu_s}$ \struck{\ill{}} ($\nu_s=\deg_{\text{sép}}T_s/k(s)$).

\uncertain{On} \uncertain{rappelle} \uncertain{le} \uncertain{raisonnement} \uncertain{du} \uncertain{cor.}\ 1,
[\uncertain{analogue} \uncertain{cor.}\ 3 \uncertain{et} \uncertain{cor.}\ 1 \ill{} \ill{} \ill{}
\uncertain{immédiat}]. Je \uncertain{veux} \uncertain{le} \uncertain{prouver} \uncertain{pour} \uncertain{le}
\uncertain{cas} \ill{} \uncertain{exact} : $\lambda^{\nu_s}$ \uncertain{points} \uncertain{dans}
\ill{} $\lambda$ \ill{} \ill{} \uncertain{fibre} \uncertain{de} $f_{*}(Y/T)$,
\uncertain{lorsque} \ill{} \ill{} \uncertain{ici} \uncertain{seulement} (\uncertain{sous} \uncertain{les}
\uncertain{conditions} \struck{\ill{} \ill{} \ill{}}).
\marginal{\ill{} \ill{} \ill{} \ill{} \ill{} \ill{} \ill{} $\lambda$ \ill{} \ill{} $\lambda$ \ill{}}

\textbf{\uncertain{Corollaire} 3} (\struck{$T/S$ \uncertain{fini}} \add{\uncertain{du} \uncertain{cor.}\ 1}).
\uncertain{Conditions} \uncertain{équivalentes} :
\begin{enumerate}
  \item[(i)] $\nu(s)\to$ \add{\uncertain{loc.}} \uncertain{constant} ;
  \item[(ii)] $X\mapsto f_{*}(X/T)$ \uncertain{transforme} \uncertain{morphismes} \uncertain{étales}
  \struck{\uncertain{finis}} \add{\uncertain{étales}} \uncertain{finis} \uncertain{en} \uncertain{morphismes}
  \uncertain{étales} \uncertain{finis} ;
  \item[(iii)] \uncertain{si} $T'$ \uncertain{est} \uncertain{étale} \uncertain{fini} \uncertain{sur} $T$,
  \uncertain{alors} $f_{*}(T'/T)$ \uncertain{est} \uncertain{étale} \uncertain{fini} ;
  \item[(iv)] $f_{*}(T\amalg T)=$ \struck{$\underline{\operatorname{Hom}}_S(T,S'\amalg S)$}
  \uncertain{est} \uncertain{étale} \uncertain{fini} \uncertain{sur} $S$.
\end{enumerate}
\note{dans la marge, des flèches relient (i), (ii), (iii), (iv) : (i) $\Rightarrow$ (ii) $\Rightarrow$ (iii)
$\Rightarrow$ (iv) et un arc de retour (iv) $\Rightarrow$ (i). La marge gauche porte aussi une note écrite en
travers, \ill{}.}

\page{54}
\uncertain{On} \uncertain{voit} \uncertain{donc} \uncertain{que} \uncertain{le} \uncertain{morphisme}
\struck{\ill{}} \uncertain{fini} $T\to S$ \uncertain{est} \uncertain{tel} \uncertain{que} $\nu(s)$
\uncertain{est} \uncertain{loc.\ constant} \add{\uncertain{dans} \uncertain{le} \uncertain{cas}}
\uncertain{si} \uncertain{c'est} \uncertain{un} \uncertain{morphisme} \uncertain{fini}
\uncertain{qui} \uncertain{est} \uncertain{composé} \uncertain{d'un} \uncertain{morphisme}
\uncertain{étale} \uncertain{fini} \uncertain{et} \uncertain{d'un} \uncertain{morphisme}
\uncertain{radiciel} [\uncertain{dans} \uncertain{le} \uncertain{cas} \uncertain{des} \uncertain{variétés}]
\uncertain{induisant} \uncertain{par} \uncertain{composition} \uncertain{des} \uncertain{morphismes},
\uncertain{extension} \uncertain{de} \uncertain{la} \uncertain{base} [\uncertain{triviale} \uncertain{sur} (i), \uncertain{par}
\uncertain{la} \uncertain{définition}].

\struck{\uncertain{Corollaire}}

\struck{\textbf{Proposition.} \uncertain{Soient} \ill{} $T,S,X,Y$ \uncertain{comme}
\uncertain{dans} \ill{} \uncertain{Prop.\ précédente}. \uncertain{Supposons}}
\struck{\uncertain{que} \ill{} \uncertain{génériquement}. \uncertain{Alors} \uncertain{pour} \uncertain{que} $X$}
\struck{\uncertain{de} $f$ \uncertain{soit} \uncertain{fibrés} \uncertain{principaux} \uncertain{de}}
\note{ce début de proposition est barré de traits obliques et enfermé dans un contour ; la dernière ligne
(« fibrés principaux de … ») est en plus rayée horizontalement.}

\textbf{Cor.\ 4.} \uncertain{Supposons} \uncertain{les} \uncertain{conditions} \uncertain{de} (i) \uncertain{remplies},
$T/S$ \uncertain{fini}, \uncertain{et} \struck{\uncertain{que} $f$ \uncertain{soit}} \uncertain{si} $g$ \uncertain{est}
\uncertain{surjectif} \struck{\uncertain{fini} \uncertain{et} \ill{} \ill{}}.
\uncertain{Alors} $f_{*}(g)$ \uncertain{est} \uncertain{surjectif}.

\textbf{Proposition.} \uncertain{Le} \uncertain{foncteur} $f_{*}$ \uncertain{transforme}
\uncertain{fibrés} \uncertain{principaux} \uncertain{sur} \uncertain{un} $Y/T$ \add{\uncertain{de} \uncertain{groupe}
\uncertain{d'une} \uncertain{certaine} \uncertain{restriction} $G$ \uncertain{de} \uncertain{constante}}
\ill{} \ill{} \uncertain{en} \uncertain{fibrés} \uncertain{principaux}.
\uncertain{Je} \uncertain{dis} \uncertain{que} \uncertain{les} \uncertain{fibrés} \uncertain{principaux}
\ill{} \ill{} \uncertain{fibrés} \uncertain{principaux} \ill{}
\uncertain{simplement} \uncertain{en} \uncertain{fibrés} \uncertain{principaux} \ill{}
\uncertain{homogènes}, \uncertain{à} \uncertain{condition} \uncertain{que} $T/S$ \uncertain{soit} \uncertain{fini}.
\uncertain{Il} \uncertain{transforme} \ill{} \struck{\ill{}} \uncertain{fibrés} \uncertain{principaux} \ill{}
\uncertain{groupes} \ill{} \ill{} \uncertain{fermés} \uncertain{séparables} \ill{} \struck{\ill{}} \ill{}
\marginal{\uncertain{Prendre} \uncertain{plutôt} \uncertain{la} \uncertain{notion} \uncertain{de} \uncertain{fibré}
\uncertain{principal} \uncertain{homogène}}
\note{la phrase se poursuit à la p. 55.}

\page{55}
\uncertain{fibrés} \uncertain{de} $=$ \uncertain{groupe}, \uncertain{pourvu} \uncertain{que}
\uncertain{l'extension} \uncertain{sur} $S$, \uncertain{sans} \uncertain{nécessairement} \uncertain{que}
$T\to S$ \uncertain{soit} \uncertain{homogène}.

\uncertain{En} \uncertain{particulier}, \uncertain{si} $G$ \uncertain{est} \uncertain{un} \uncertain{groupe}
\uncertain{fini} \uncertain{ordinaire}, \uncertain{un} \struck{\ill{}} \uncertain{revêtement}
\struck{\ill{}} \uncertain{principal} \uncertain{de} \uncertain{groupe} $G$ \uncertain{de} $Y/T$
\uncertain{est} \uncertain{transformé} \uncertain{en} \uncertain{un} \struck{\ill{}} \uncertain{fibré}
\uncertain{principal} \uncertain{homogène} \uncertain{sur} $f_{*}(Y/T)$, \uncertain{de} \uncertain{groupe}
\[
  f_{*}(Y/T)\times_S f_{*}(T\times G/T).
\]
\note{dans la dernière formule, un mot biffé et illisible précède $T\times G$.}

\uncertain{Cela} \uncertain{conduit} \uncertain{à} \uncertain{l'intéressante} \uncertain{question}
\uncertain{de} \uncertain{déterminer}, \uncertain{pour} \uncertain{un} \uncertain{ens.}\ \uncertain{fini} $I$
\struck{\uncertain{ordinaire}}, \uncertain{la} \uncertain{variété} \add{\uncertain{étale}}
\[
  f_{*}(T\times I/T)=\varphi(I)
\]
\struck{\uncertain{qui} \uncertain{est} \uncertain{telle} \uncertain{que}}
\struck{\uncertain{donc} $\varphi(T)(I)$ \uncertain{est} \uncertain{un} \uncertain{foncteur}}
\struck{\uncertain{covariant} \uncertain{en} $I$ / \uncertain{qui} \uncertain{ne} \uncertain{dépend} \uncertain{que}}
\struck{\uncertain{de} \uncertain{la} \uncertain{nature} \uncertain{étale} \uncertain{de} $S$. \uncertain{Plaçons}}
\struck{\uncertain{nous} \uncertain{dans} \uncertain{le} \uncertain{cas} \uncertain{où} \uncertain{on} \uncertain{a}
\uncertain{sur} $S$, \uncertain{supposé}}
\struck{\uncertain{connexe}, \uncertain{un} \uncertain{point} \uncertain{géométrique}
\uncertain{fixé} $a$, \uncertain{de} \uncertain{sorte} \uncertain{que}}
\struck{\uncertain{le} \uncertain{foncteur} $I\mapsto\varphi_a(T)(I)$ (\uncertain{Ens} \uncertain{finis})
\uncertain{dépend}}
\struck{\uncertain{de} \uncertain{la} \uncertain{donnée} \uncertain{du} \uncertain{groupe} $\pi_1(S,a)$
\uncertain{opérant}. \uncertain{Et} \uncertain{étant} \uncertain{donné}}
\struck{\uncertain{un} \uncertain{foncteur} \uncertain{à} \uncertain{valeurs} \uncertain{dans} \uncertain{les}
\uncertain{ens.\ finis},}
\struck{\uncertain{il} \uncertain{est} \uncertain{de} \uncertain{la} \uncertain{forme}
$\operatorname{Hom}_{\text{Ens.f}}(A_a(T),I)$,}
\struck{[\uncertain{petit} \uncertain{raisonnement} \ill{}], \uncertain{et} \uncertain{comme}}
\struck{$\pi_1(S,a)$ \uncertain{opère} \uncertain{dans} \uncertain{ce} \uncertain{foncteur}, \uncertain{il}
\uncertain{opère}}
\struck{\uncertain{dans} $A_a(T)$, \uncertain{qui} \uncertain{est} \uncertain{un} \uncertain{ainsi} \uncertain{défini}.}
\note{depuis « qui est telle que », tout le bas de la page est barré de deux longues diagonales ; une note
marginale oblique, à gauche, renvoie par un trait à ce passage : \ill{}. Le $\varphi_a(T)(I)$ porte un $I$ surchargé.}

\page{56}
\struck{\uncertain{de} $f$ \uncertain{sur} \uncertain{naturelle} \uncertain{un} \uncertain{revêtement}
\uncertain{étale}}
\struck{[$T_{\text{sép}}$ \ill{} \ill{} \ill{} $S$, \uncertain{rapporté} \ill{}]}
\struck{$\underline{A}(T)$ \uncertain{de}} $f^{\text{ét}}_{*}(T)$, \uncertain{qui} \uncertain{dépend}
\uncertain{fonctoriellement}
\uncertain{de} $T$, \struck{\uncertain{et} \uncertain{transforme}} \uncertain{le} \uncertain{foncteur}
$f^{\text{ét}}_{*}$ \uncertain{en} \uncertain{fonctions}
\uncertain{d'ailleurs} \ill{} \ill{} \ill{} \struck{\uncertain{produits}} \uncertain{sommes}.
\uncertain{On} \uncertain{va} \uncertain{définir} \uncertain{un} $S$-\uncertain{morphisme}
\[
  T\to\overline{T}
\]
i.e.\ \uncertain{un} \uncertain{élément} \uncertain{de} $\underline{\operatorname{Hom}}_S(T,\overline{T})$, i.e.\ \uncertain{une}
\uncertain{section} \uncertain{de} $\underline{\operatorname{Hom}}_S(T,T)$.
\note{depuis le début de la page jusqu'ici, le texte est barré d'une longue diagonale et cerné d'un trait ;
seuls « $f^{\text{ét}}_{*}(T)$ » et la suite semblent avoir été gardés après la rature de l'en-tête encadré.
Le partage entre biffé et conservé est proposé, non assuré.}

\uncertain{Plus} \uncertain{généralement}, \uncertain{si} $T$ \uncertain{est} \uncertain{donné}, $S'\mapsto
\underline{\operatorname{Hom}}_S(T,S')$ \uncertain{est} \uncertain{un} \uncertain{foncteur} \uncertain{des}
\uncertain{revêtements} \uncertain{étales} \uncertain{de} $S$ \uncertain{dans} \uncertain{les} \uncertain{rev.}\
\uncertain{étales} \uncertain{de} $S$, \uncertain{exact} \uncertain{à} \uncertain{gauche}.
\struck{\uncertain{Si} \uncertain{l'on} \uncertain{suppose} \uncertain{seulement} \uncertain{donné} \uncertain{un}}
\struck{\uncertain{foncteur} \uncertain{de} $S$, \uncertain{supposé} \uncertain{connexe} \ill{}}
\struck{\ill{} \ill{} \ill{} \uncertain{foncteur}}
\struck{\uncertain{déterminé} \uncertain{par} \uncertain{un} \uncertain{élément} \uncertain{fini} \uncertain{à}}
\struck{\uncertain{prouver} \uncertain{que} \uncertain{le} \uncertain{foncteur} \uncertain{est} \uncertain{de} \uncertain{la}}
\struck{\uncertain{forme} \uncertain{lim} $\underline{\operatorname{Hom}}($}
\uncertain{Fixons} \uncertain{un}
\struck{\uncertain{point} \uncertain{géométrique}} \uncertain{de} $S$ \uncertain{supposé}
\uncertain{connexe} (\uncertain{ce} \uncertain{qu'on} \uncertain{devrait} \uncertain{pas} \uncertain{nécessairement})
\struck{\uncertain{Alors} \uncertain{le} \uncertain{foncteur}} \uncertain{la} \uncertain{donnée}
\marginal{(\uncertain{i.e.}\ \uncertain{la} \uncertain{fonctorialité} (\ill{}) \uncertain{suppose} \uncertain{un}
\uncertain{élément} \uncertain{du} \uncertain{foncteur} $S'\mapsto\underline{\operatorname{Hom}}_S(T,S')$)}
\note{le bas de la page est un enchevêtrement de ratures (traits horizontaux, diagonales, boucles) ; la note
marginale oblique est rattachée par un trait pointillé à la ligne « Fixons ». La phrase continue p. 57.}

\page{57}
\uncertain{d'un} \uncertain{foncteur} \uncertain{Ens} $\varphi_{T,a}(E)$ \uncertain{des} $\pi_1(S,a)$-\uncertain{ensembles}
\uncertain{finis} \uncertain{en}, \uncertain{ensembles} \uncertain{finis}, \uncertain{foncteur} \uncertain{qui} \uncertain{est}
\uncertain{exact} \uncertain{à} \uncertain{gauche}. \uncertain{Je} \uncertain{dis} \uncertain{qu'un} \uncertain{tel}
\uncertain{foncteur} \uncertain{est} \uncertain{nécessairement} \uncertain{de} \uncertain{la}
\uncertain{forme} $\underline{\operatorname{Hom}}_{\pi_1}(A,E)$, \uncertain{où} $A$ \uncertain{est}
\struck{\ill{}}
\uncertain{un} \uncertain{ens.}, \uncertain{dépendant} \uncertain{fonctoriellement} \uncertain{du}
\uncertain{foncteur} $\varphi_{T,a}$ \uncertain{qui} \uncertain{lui} \uncertain{donne} \uncertain{naissance}
[\uncertain{il} \uncertain{suffit} \uncertain{d'avoir} \uncertain{fonctoriellement} \uncertain{en} \uncertain{fonction}
\uncertain{de} \uncertain{finis} \uncertain{au} \uncertain{lieu} \uncertain{de} \uncertain{finis} \uncertain{en} \uncertain{nombre}]. \uncertain{On}
\uncertain{posera} $A=A(T)$, \struck{\uncertain{On}} \uncertain{et} \uncertain{on} \uncertain{dénotera}
\uncertain{par} $f^{\text{ét}}_{*}(T)$ \uncertain{le} \uncertain{revêtement} \struck{\uncertain{équivalent}}
\add{\uncertain{étale}} \uncertain{de} $S$ \uncertain{correspondant}. \uncertain{On} \uncertain{a} \uncertain{donc}
\uncertain{un} \uncertain{isom.\ de} \uncertain{bifoncteurs}
\[
  (*)\qquad \underline{\operatorname{Hom}}_S\bigl(f^{\text{ét}}_{*}(T),S'\bigr)
  =\underline{\operatorname{Hom}}_S(T,S')
\]
[\uncertain{sur} $\operatorname{Rev\,\acute{e}t}(T)\times\operatorname{Rev\,\acute{e}t}(S)$].
\struck{\uncertain{D'autre} \uncertain{part} \uncertain{pour} \uncertain{tout} $T$, \uncertain{le} \uncertain{foncteur}}
\struck{\uncertain{en} $S'$ \uncertain{est} \uncertain{représentable}. \uncertain{Prenons} \uncertain{donc}
$S'=T$, \uncertain{et}}
\note{ces deux lignes sont encadrées et barrées de traits obliques.}
\struck{\uncertain{Le} \uncertain{morphisme} \ill{}} \uncertain{les} \uncertain{notations}, \uncertain{on} \uncertain{dénote}
\uncertain{tous}, \uncertain{on} \uncertain{trouve} \uncertain{un}
\[
  \underline{\operatorname{Hom}}_S\bigl(f^{\text{ét}}_{*}(T),S'\bigr)\simeq\underline{\operatorname{Hom}}_S(T,S')
\]
\uncertain{d'où}, \uncertain{en} \uncertain{faisant} $S'=f^{\text{ét}}_{*}(T)$, \uncertain{un} \uncertain{homomorphisme}
\marginal{\uncertain{à} \uncertain{distinguer} \uncertain{de} \uncertain{l'intérieur} \ill{} \ill{} \uncertain{de} \uncertain{ces}
\uncertain{objets}, \uncertain{qu'il} \uncertain{faut} \uncertain{établir}, \uncertain{pour} \uncertain{un} \ill{} \ill{}
\uncertain{précédemment} (*) \ill{} \ill{} \ill{} \uncertain{définie}}
\note{la note marginale, écrite en oblique et appelée par (*), longe toute la partie gauche de la page ; la
phrase continue p. 58.}

\page{58}
\uncertain{canonique} :
\[
  T\longrightarrow f^{\text{ét}}_{*}(T)
\]
\uncertain{tel} \uncertain{que} \uncertain{pour} \uncertain{tout} $S'$, \uncertain{l'homomorphisme}
\uncertain{correspondant}
\[
  \underline{\operatorname{Hom}}_S\bigl(f^{\text{ét}}_{*}(T),S'\bigr)\longrightarrow\underline{\operatorname{Hom}}_S(T,S')
\]
\struck{\ill{}} \uncertain{soit} \uncertain{un} \uncertain{isom.\ pour} \uncertain{tout} $S'$ (\uncertain{Précisions} …).
\uncertain{Il} \uncertain{en} \uncertain{résulte} \uncertain{que} \uncertain{pour} \uncertain{tout} \struck{$s\in S$,}
\uncertain{changement} \uncertain{de} \uncertain{base}
\struck{$\underline{\operatorname{Hom}}_{k(s)}(\overline{T}_s,S'_s)\to$}
[\uncertain{base} (\uncertain{pt} \uncertain{géom.}\ $k(s)$) $S$ : $\operatorname{Spec}(k(\bar s))$]
\uncertain{la} \uncertain{construction} \uncertain{commute} \uncertain{aux} \uncertain{changements} \uncertain{de} \uncertain{base}…

\uncertain{En} \uncertain{faisant} $S'=S\times I$ ($I$ \uncertain{ens.\ fini}) \uncertain{on} \uncertain{voit}
\uncertain{d'abord} \uncertain{que}
\[
  T\to f^{\text{ét}}_{*}(T)
\]
\uncertain{est} \uncertain{surjectif} \uncertain{et} \add{(\uncertain{si} $T/S$ \uncertain{est} \uncertain{fini})}
\uncertain{à} \uncertain{fibres} \uncertain{géométriques} \uncertain{connexes}
(\uncertain{on} \uncertain{considère} \add{\uncertain{les} \uncertain{fibres} \uncertain{des}}
\uncertain{fibres} \uncertain{géom.\ pour} $\bar s\in S$).
\uncertain{Il} \uncertain{en} \uncertain{résulte}, \uncertain{si} $T\to\overline{T}$ \uncertain{est} \uncertain{un}
\uncertain{morphisme} \uncertain{surjectif} \uncertain{et} \uncertain{fini} \uncertain{injectif} \uncertain{de}
$T$ (\uncertain{fini} \add{\uncertain{étale}} \uncertain{sur} $S$) \uncertain{sur} $\overline{T}$ \uncertain{rev.\ étale}
\uncertain{de} $S$,
\uncertain{alors} \uncertain{on} \uncertain{a}
\[
  \underline{\operatorname{Hom}}_S(\overline{T},S')=\underline{\operatorname{Hom}}_S(T,S')
\]
\uncertain{pour} \uncertain{tout} $S'$ \uncertain{rev.} \uncertain{étale} \uncertain{fini} \uncertain{sur} $S$
[\uncertain{démonstration} \ill{}] \uncertain{et} \uncertain{par} \uncertain{conséquent} \uncertain{le}
$\overline{T}$ \uncertain{dont} \uncertain{il} \uncertain{est} \uncertain{question} \uncertain{se}
\uncertain{trouve} \uncertain{vrai} \uncertain{pour} \uncertain{cat.\ de} \uncertain{la} \uncertain{base}, \add{$(S)$}
\uncertain{ce} \uncertain{qui} \uncertain{donne} \struck{\ill{}}
\[
  \underline{\operatorname{Hom}}_S(\overline{T},S')\simeq\underline{\operatorname{Hom}}_S(T,S'),
\]
i.e.\ $\overline{T}=f^{\text{ét}}_{*}(T)$.
\note{le « $=$ » et le « $\simeq$ » des deux formules sont lus tels quels ; la seconde semble porter une double
flèche ($\rightrightarrows$) plutôt qu'un $\simeq$, lecture incertaine.}

\page{60}
\section{(1) \uncertain{Sit} \uncertain{up} \uncertain{optimal}}
\note{titre souligné, en haut de la page, précédé de « (1) » ; lecture très incertaine. La page ouvre une
nouvelle rédaction, qui se poursuit au-delà de ce lot.}

\uncertain{Soit} $S$ \uncertain{un} \uncertain{préschéma}. \struck{\ill{}}
\uncertain{Soit} $\underline{W}$ \uncertain{un} \uncertain{foncteur} \uncertain{covariant} \uncertain{de} \uncertain{la}
\uncertain{catégorie} \uncertain{des} \struck{$S$-}\uncertain{préschémas} \uncertain{dans} \uncertain{la}
\uncertain{catégorie} \uncertain{des} \struck{$E$-}\uncertain{préschémas}.
\uncertain{On} \uncertain{suppose} $\underline{W}$ « \uncertain{de} \uncertain{nature} \uncertain{locale} », i.e.
\begin{enumerate}
  \item[(i)] \uncertain{Si} $i\colon T\to T'$ \uncertain{est} \uncertain{une} \uncertain{immersion} \uncertain{ouverte},
  \uncertain{alors} $\underline{W}(i)\colon\underline{W}(T)\to\underline{W}(T')$ \uncertain{est} \uncertain{une}
  \uncertain{immersion} \uncertain{ouverte}. [\uncertain{Donc} \uncertain{si} $U$ \uncertain{ouvert} \uncertain{dans}
  $T$, \uncertain{on} \uncertain{identifie} $\underline{W}(U)$ \uncertain{à} \uncertain{un} \uncertain{ouvert}
  \uncertain{de} $\underline{W}(T)$] ;
  \item[(ii)] $\underline{W}(\bigcup U_i)=\bigcup_i\underline{W}(U_i)$ \uncertain{pour} \uncertain{une}
  \uncertain{famille} \uncertain{d'ouverts} \uncertain{de} $T$ ;
  \item[(iii)] $\underline{W}(U\cap V)=\underline{W}(U)\cap\underline{W}(V)$.
\end{enumerate}
\marginal{\ill{} \uncertain{infinitésimale} \ill{} \ill{}}
\note{la note marginale oblique est rattachée par une accolade aux conditions (ii)--(iii).}

\textbf{Exemple I.} \uncertain{Soit} $S'$ \uncertain{un} $S$-\uncertain{préschéma}.
\uncertain{Posons} $\underline{W}_{S'}(T)=T\times_S S'$ ; \uncertain{on} \uncertain{obtient} \uncertain{ainsi}
\uncertain{un} \uncertain{foncteur} \uncertain{covariant} \uncertain{de} \uncertain{nature} \uncertain{locale}
\uncertain{des} $S$-\uncertain{préschémas} \uncertain{dans} \uncertain{les} \add{$S'$}-\uncertain{préschémas}
(\uncertain{le} \uncertain{foncteur} \uncertain{de} \uncertain{changement} \uncertain{de} \uncertain{base}).
\uncertain{Cas} I$'$ : $S'$ \uncertain{est} \uncertain{fini} \uncertain{et} \uncertain{loc.\ libre} \uncertain{sur} $S$.

\textbf{Exemple II.} \uncertain{Soit} $W$ \uncertain{un} $A$-\uncertain{préschéma} \uncertain{en}
\ill{} \uncertain{commutatifs} \uncertain{sur} $S$, \uncertain{supposé} \uncertain{plat}
\uncertain{pour} \uncertain{tout} $T/S$, \uncertain{le} \uncertain{faisceau} \uncertain{d'anneaux}
$\Gamma(W_T/T)$ \uncertain{sur} $T$ \uncertain{forme} \uncertain{de} $T$ \uncertain{un}
\uncertain{préschéma} (\uncertain{cf} \uncertain{exemple} \uncertain{plus} \uncertain{bas}).
\uncertain{On} \uncertain{a} \uncertain{ainsi}
\marginal{\uncertain{annulé} \uncertain{en} \uncertain{anneaux} \uncertain{locaux}}
\note{la note marginale, en bas à gauche, renvoie par un cercle au mot « préschéma ». La phrase « On a ainsi »
est coupée en bas de page ; l'argument continue au-delà de ce lot.}

\end{document}
